\chapter{几何与波}

\providecommand{\tr}{\operatorname{tr}}
\providecommand{\by}{\mathbf{y}}
\providecommand{\bb}{\mathbf{b}}

本章讨论正交宇宙（\emph{Orthogonal}）中的波与对称性。我们首先从理想化的平面波出发，说明在黎曼几何（Riemannian geometry）中"光"的传播矢量不再是零矢量，其空间频率 $\kappa$ 与时间频率 $\nu$ 满足 $\kappa^{2}+\nu^{2}=\nu_{\max}^{2}$；随后导出黎曼标量波方程（RSW）与黎曼矢量波方程（RVW）。进阶部分系统介绍对称群的表示论——从群表示的基本概念、可约与不可约表示、酉表示，到 SO(3)、SO(4) 的双覆盖以及 SU(2) 的自旋-$j$ 表示——为理解正交宇宙中粒子与波所服从的对称性奠定数学基础。

\section{平面波}

考虑一束穿过我们自己所处的洛伦兹宇宙的理想化光波：一根无限长、频率精确定义的正弦波。按惯例把这个频率记作 $\nu$。暂时忽略偏振问题，只用一个数 $A$ 来描述波在任意地点与时刻的强度。

设波沿坐标系的 $x$ 轴传播。在光速取为 1 的单位制下，这列波可表示为

\[
A(x,y,z,t)=A_{0}\sin\!\bigl(2\pi\nu\,(x-t)\bigr)
\]

\orthofig{images/03__001.png}{我们宇宙中的光波}

频率的倒数称为波的\textbf{周期}（period）$\tau$，即完成一个完整周期所需的时间。由于正弦函数在其宗量增加或减少 $2\pi$ 时精确重复，波 $A(x,y,z,t)$ 会在 $t$ 增加或减少 $\tau=1/\nu$ 时重复：

\[
A(x,y,z,t\pm\tau)=A(x,y,z,t)
\]

在光速取 1 的单位制下，光的\textbf{波长}（wavelength）与周期数值相等：$\lambda=1/\nu$。波 $A$ 会在 $x$ 增加或减少一个波长时重复：

\[
A(x\pm\lambda,y,z,t)=A(x,y,z,t)
\]

现在定义一个零矢量（null vector）$\bk$，它在我们坐标系中的分量为

\[
k^{x}=2\pi\nu,\qquad k^{y}=0,\qquad k^{z}=0,\qquad k^{t}=2\pi\nu
\]

我们称它为光波的\textbf{传播矢量}（propagation vector）。于是波的振幅可改写为

\[
A(x,y,z,t)=A_{0}\sin(k^{x}x+k^{y}y+k^{z}z-k^{t}t)
\]

表达式 $k^{x}x+k^{y}y+k^{z}z-k^{t}t$ 是把传播矢量 $\bk$ 的各分量与 $\bx=(x,y,z,t)$ 的相应分量相乘后再相加，这让人联想到黎曼几何中定义的点积，唯一的差别是 $t$ 分量乘积前带了一个负号。事实上它正是 $\bk\cdot\bx$ 的洛伦兹版本，并且是我们的洛伦兹宇宙中所有观测者——无论怎样运动——都会一致同意的量。

\orthofig{images/03__004.png}{红外极限}
\orthofig{images/03__002.png}{黎曼宇宙中的慢光波}
\orthofig{images/03__003.png}{黎曼宇宙中的快光波}
\orthofig{images/03__005.png}{紫外极限}

这提示我们，在\emph{黎曼}宇宙真空中传播的同类波应当取形式

\[
A(x,y,z,t)=A_{0}\sin(\bk\cdot\bx)
\]

其中点积是黎曼点积。在黎曼几何中\emph{不存在零矢量}：四维空间中的所有方向本质上都一样。因此表面上看，$\bk$ 可以是任意矢量。

不过我们总可以选取坐标系使 $\bk$ 落在 $xt$ 平面内，于是 $k^{y}=k^{z}=0$，从而

\[
A(x,y,z,t)=A_{0}\sin(\bk\cdot\bx)=A_{0}\sin(k^{x}x+k^{t}t)
\]

这列波的\textbf{波前}（wavefront）——即"波峰"——出现在 $\sin(\bk\cdot\bx)=1$ 处，也就是 $\bk\cdot\bx=\pi/2+2n\pi$（$n$ 为整数）处。若找到一点使 $\bk\cdot\bx$ 取这样的值，则在四维空间中存在三个与 $\bk$ 垂直的方向，沿它们移动不会改变 $\bk\cdot\bx$，所以波前是一系列等间距的三维区域，且这些区域与 $\bk$ 垂直。在我们的图中只画出四维空间的一张二维切片，于是这些区域表现为一列平行线。

这列波的波长 $\lambda$ 与周期 $\tau$ 是多少？不难看出

\[
A\!\left(x\pm\frac{2\pi}{k^{x}},y,z,t\right)=A(x,y,z,t),\qquad
A\!\left(x,y,z,t\pm\frac{2\pi}{k^{t}}\right)=A(x,y,z,t)
\]

所以波长与周期为

\[
\lambda=\frac{2\pi}{k^{x}},\qquad \tau=\frac{2\pi}{k^{t}}
\]

我们称波长的倒数为\emph{空间频率}（spatial frequency）$\kappa$；它完全类比于时间频率 $\nu$，告诉我们单位距离内可容纳多少个波周期。于是传播矢量 $\bk$ 的分量可以写成

\[
k^{x}=\frac{2\pi}{\lambda}=2\pi\kappa,\qquad k^{y}=0,\qquad k^{z}=0,\qquad k^{t}=\frac{2\pi}{\tau}=2\pi\nu
\]

并且空间频率 $\kappa$、时间频率 $\nu$ 与传播矢量大小 $|\bk|$ 之间满足

\[
|\bk|^{2}=(k^{x})^{2}+(k^{t})^{2}=(2\pi\kappa)^{2}+(2\pi\nu)^{2}
\]

于是

\[
\kappa^{2}+\nu^{2}=\frac{|\bk|^{2}}{4\pi^{2}}
\]

黎曼宇宙中不同速度的观测者都会同意 $|\bk|$（传播矢量大小）的取值，正如洛伦兹宇宙中的不同观测者都会同意与光波相伴的传播矢量是零矢量一样。因此我们假设，黎曼宇宙中相当于"光"的那种物理现象，由一批共享同一 $|\bk|$ 值的波组成。

固定 $|\bk|$ 的几何效果，是固定波前之间的间距——它不是在某一观测者所选定的坐标方向上量得的，而是垂直于波前、沿传播矢量方向量得的。黎曼空间中这个间距恒等于 $2\pi/|\bk|$。

〔我们稍后会看到，在现象的量子力学版本中，$|\bk|$ 的取值与\emph{这些波所对应粒子的静止质量}有关，所以"所有光具有相同 $|\bk|$"其实只是在坚持我们选定某个粒子来扮演我们宇宙中光子的角色。光子没有静止质量，但在黎曼宇宙中既然没有零矢量，也就不存在无静止质量的粒子，因此我们必须选取\emph{某个}非零值。〕

一旦 $|\bk|$ 固定，就可见随着空间频率 $\kappa$ 变大，时间频率 $\nu$ 必然变小，反之亦然，以使其平方和保持恒定。若两个频率之一为零，另一个就达到最大可能值，记作 $\nu_{\max}$，它与 $|\bk|$ 成正比：

\[
\nu_{\max}=\frac{|\bk|}{2\pi}
\]

于是 $\kappa$ 与 $\nu$ 的关系变为

\[
\kappa^{2}+\nu^{2}=\nu_{\max}^{2}
\]

正如我们在讨论双勾股定理时所见，这正是沿两根垂直轴测量的空间频率之间应有的关系。在黎曼宇宙中，时间频率与沿任何其他维度测得的频率没有区别，所以数学完全相同。频率 $\nu_{\max}$ 就是以最直接的方式、沿四维空间中垂直于波前的方向、单位距离内的波周期数。

\orthofig{images/03__006.png}{波长等量与速度的关系}

现在假设一个光脉冲的世界线沿传播矢量延伸，则光的速度 $v$ 等于 $\bk$ 的 $x$ 与 $t$ 分量之比：

\[
v=\frac{k^{x}}{k^{t}}=\frac{\kappa}{\nu}
\]

最慢的光 $v=0$，其波长无限长 $\lambda=\infty$，我们称之为"红外极限"，其时间频率为 $\nu_{\max}$。

最快的光 $v=\infty$，其波长最小 $\lambda=1/\nu_{\max}$，我们称之为"紫外极限"，其时间频率为零，因而周期无限长。

\orthofig{images/03__007.png}{平面波}

但我们凭什么假设光脉冲的世界线沿传播矢量延伸？

到目前为止所描述的波都是\textbf{平面波}（plane wave）：它们的波前是空间中的无限平面，也就是四维空间中无限超平面的切片。它们数学上非常简单，却显然高度理想化——而平面波的振幅当然\emph{不会}集中成任何类似平行于传播矢量的世界线的东西。右图中振幅最大的亮带\emph{垂直}于传播矢量，并且随着时间推移，这些波前沿与传播矢量\emph{相反}的方向移动！

但这些是占据整个四维空间的无限波。若要得到一个\emph{局域化}的波，就需要把不同频率的波叠加起来。这些波相互加强的位置不会以任何单一成分波的波前速度移动；它的速度是系统的群速度，而不是波前的相速度。

我们可以通过考察使两列频率略有差异的平面波保持同相所需的 $x$ 与 $t$ 之间的关系，来求波叠加后的群速度。设第一列波的空间、时间频率为 $\kappa,\nu$，第二列为 $\kappa+\Delta\kappa,\nu+\Delta\nu$，两列波相位保持相同要求

\[
\kappa x+\nu t=(\kappa+\Delta\kappa)x+(\nu+\Delta\nu)t
\]

即

\[
x=-\frac{\Delta\nu}{\Delta\kappa}\,t
\]

所以 $v_{\text{group}}=-\dd\nu/\dd\kappa$。把它用到 $\nu$ 与 $\kappa$ 的关系上：

\[
\nu=\sqrt{\nu_{\max}^{2}-\kappa^{2}},\qquad
v_{\text{group}}=-\frac{\dd\nu}{\dd\kappa}=\frac{\kappa}{\sqrt{\nu_{\max}^{2}-\kappa^{2}}}=\frac{\kappa}{\nu}
\]

这与我们直接从传播矢量得到的 $k^{x}/k^{t}$ 一致。

\orthofig{images/03__008.png}{略微不同的平面波之和}

左图由总共 61 列平面波叠加而成（并非所有标记传播矢量的箭头都足够明亮可见），它给出了黎曼宇宙中一个光脉冲历史的合理图像。更真实的情形是叠加一个连续的平面波谱，而非有限个波。

叠加结果大致平行于平均传播矢量，但它当然不是一条完美的锐利世界线，也不是等宽的"世界管"或"世界带"。不过脉冲随时间逐渐展宽这一事实完全在意料之中：黎曼宇宙中不同频率的光以不同速度传播，所以任何由许多不同频率构成的局域波都会如此逐渐色散。

\section{标量波}

上一节讨论了几种非常简单的波的数学表达式，但我们真正想要的是黎曼宇宙真空中\emph{一切}波——无论形状如何——都满足的方程。

我们自己宇宙中的波满足\textbf{偏微分方程}（partial differential equation），即关于波的强度沿各方向变化率之间的某种关系。若某量 $A$ 依赖于若干变量，例如我们在四维空间某个区域上放置的坐标 $x,y,z,t$，则 $A$ 对其中一个变量的\textbf{偏导数}（partial derivative），就是当只改变该变量、保持其余变量不变时 $A$ 的\emph{变化率}。

例如，考虑传播矢量为 $\bk$ 的黎曼平面波

\[
A=A_{0}\sin(\bk\cdot\bx)=A_{0}\sin(k^{x}x+k^{y}y+k^{z}z+k^{t}t)
\]

若固定 $y,z,t$ 而改变 $x$，$A$ 将按正弦波动。正弦函数对宗量的变化率就是同一宗量的余弦，但由于这里正弦函数的宗量中 $x$ 乘了因子 $k^{x}$，$A$ 的变化率也要乘上该因子；这只是微积分中称为链式法则的简单规则的一个例子。

$A$ 对 $x$ 的偏导数记作 $\partial_{x}A$。这个记号初次使用可能略显唬人，但其含义其实非常简单：设想空间中一条线，线上除 $x$ 外所有坐标都固定，然后考虑所要求导的量如何\emph{沿该线}变化，忽略它在别处的行为。用此记号可写

\[
\partial_{x}A=k^{x}A_{0}\cos(k^{x}x+k^{y}y+k^{z}z+k^{t}t)
\]

如果再对 $x$ 求\emph{这个}量的变化率呢？余弦函数对宗量的变化率是同一宗量的正弦的\emph{相反数}，同时链式法则又给出一个因子 $k^{x}$。二阶偏导数记作 $\partial_{x}^{2}A$，故

\[
\partial_{x}^{2}A=-(k^{x})^{2}A_{0}\sin(k^{x}x+k^{y}y+k^{z}z+k^{t}t)=-(k^{x})^{2}A
\]

对另外三个坐标的二阶偏导数为

\[
\partial_{y}^{2}A=-(k^{y})^{2}A,\qquad \partial_{z}^{2}A=-(k^{z})^{2}A,\qquad \partial_{t}^{2}A=-(k^{t})^{2}A
\]

四者相加得

\[
\partial_{x}^{2}A+\partial_{y}^{2}A+\partial_{z}^{2}A+\partial_{t}^{2}A
=-\bigl[(k^{x})^{2}+(k^{y})^{2}+(k^{z})^{2}+(k^{t})^{2}\bigr]A=-|\bk|^{2}A
\]

这里我们做的，是用在每个坐标方向对 $A$ 求二阶变化率的操作，把波的原始函数乘以一个正比于该方向频率平方的数；然后把\emph{所有}二阶变化率相加，就得到平方和，它等于 $|\bk|^{2}$——一个完全独立于传播矢量具体方向的因子。

现在 $|\bk|=2\pi\nu_{\max}$。为简洁起见（避免到处出现 $2\pi$ 因子），我们采用"角频率"记号，用符号 $\omega$ 表示，它等于普通频率 $\nu$ 的 $2\pi$ 倍。于是 $|\bk|=2\pi\nu_{\max}=\wms$，上式可改写为

\begin{equation}
\partial_{x}^{2}A+\partial_{y}^{2}A+\partial_{z}^{2}A+\partial_{t}^{2}A+\wms^{2}A=0 \tag{RSW}
\end{equation}

这个方程不再涉及原来那个平面波的任何具体细节，它只涉及一个在四维空间上变化的量 $A$，以及一个与波形状无关的数 $\wms$。

我们称此方程为\textbf{黎曼标量波方程}（Riemannian Scalar Wave equation, RSW）。"标量"（scalar）是指在四维空间每一点处的、所有观测者都能一致同意的一个数，这与矢量分量不同——后者依赖于坐标系的选择。此方程是我们自己宇宙物理学中出现的、称为亥姆霍兹（Helmholtz）方程的四维版本。

求导运算是线性的：若 $s$ 为常数，$A$、$B$ 为 $x$ 的函数，则

\[
\partial_{x}(sA)=s\,\partial_{x}A,\qquad \partial_{x}(A+B)=\partial_{x}A+\partial_{x}B
\]

二阶变化率同样如此，乘以常数如 $\wms^{2}$ 当然也如此。所以 RSW 方程是一个\textbf{线性方程}：取两个解相加，或把一个解乘以常数，结果仍是解。既然所有共享同一 $\wms$ 的黎曼平面波都满足 RSW 方程，那么由任意多个这样的平面波叠加而成的波也满足它。例如，上一节末尾由 61 列平面波叠加而成的波就满足 RSW 方程。

RSW 方程有一个——可能是灾难性的——问题。考虑波

\[
A=\sin\!\left[\tfrac{5}{4}\wms x\right]\exp\!\left[\tfrac{3}{4}\wms t\right]
\]

指数函数的变化率仍是指数函数，而二阶变化率给出原函数的\emph{正}倍数，不同于正弦、余弦给出的负倍数。于是

\[
\partial_{x}^{2}A=-\tfrac{25}{16}\wms^{2}A,\qquad \partial_{y}^{2}A=0,\qquad \partial_{z}^{2}A=0,\qquad \partial_{t}^{2}A=\tfrac{9}{16}\wms^{2}A
\]

把这些项相加可见，上面给出的这个波满足 RSW 方程。但它描述的却是一个随时间\emph{指数增长}的波！指数因子之所以可能出现，是因为它二阶变化率的正号被另一个坐标方向上大于 $\nu_{\max}$ 的空间频率所平衡——即 $\sin[(5/4)\wms x]$ 中的角频率 $(5/4)\wms$。仅仅声称"$\nu_{\max}$ 是最大可能频率"并不足以使它成立，而单靠 RSW 方程也无法强制这一限制。但如果允许这样的波存在，那么一个频率足够高的起初极微小的扰动，将迅速增长并吞没其他一切。

有一种办法可以避免这个问题，但它的发现是小说的重要情节，所以我在此不予剧透。

\section{矢量波}

我们宇宙中的光是电磁波。电磁场不是由逐点变化的单一数值刻画，而是涉及\emph{矢量}。因此我们想弄清，在黎曼宇宙中什么样的矢量波才说得通。

设固定一个四维矢量 $\bA_{0}$，它在我们坐标系中的分量为 $A_{0}^{x},A_{0}^{y},A_{0}^{z},A_{0}^{t}$。则可考虑矢量平面波

\[
\bA(x,y,z,t)=\bA_{0}\sin(\bk\cdot\bx)
\]

它与此前考虑的标量平面波非常相似。事实上，只要 $|\bk|=\wms$，$\bA$ 在我们坐标系中的\emph{每一个分量}都满足 RSW 方程。当然这些分量在不同坐标系中会不同……但若它们在一个坐标系中满足 RSW 方程，则在任何其他坐标系中也满足。（注意这里说的是不同的\emph{直角}坐标系，而非球坐标之类。）

这足够简单，但有一个小麻烦需要处理。虽然 $\bA$ 的分量 $A^{x},A^{y},A^{z},A^{t}$ 会随坐标系改变，但从上面写出的矢量波可以构造出一个标量波：

\[
D(x,y,z,t)=\bk\cdot\bA(x,y,z,t)=(\bk\cdot\bA_{0})\sin(\bk\cdot\bx)
\]

这个波是真正的\emph{标量}波——所有观测者无论怎样选取坐标系都会同意它的取值。这看似不重要，但我们在自己宇宙中的所有物理经验都表明，标量波与矢量波将对应不同的现象。在量子力学中，标量波对应自旋为零的粒子，矢量波对应携带某一单位自旋的粒子——光子便是一例。所以我们期望黎曼宇宙中相当于光的现象应是一个纯矢量波，不应附带如 $D$ 这样的标量波。

于是，我们对像 $\bA$ 这样的矢量平面波提出的要求是 $\bk\cdot\bA_{0}=0$。施加此条件意味着描述该波自身的矢量与传播矢量正交，从而把 $\bA_{0}$ 限制在一个与 $\bk$ 正交的三维子空间内。既然该子空间是三维的，黎曼光将有\emph{三}个独立的偏振态，而我们自己宇宙中的光只有两个。等我们详细考察黎曼电磁学时还会再谈这一点。

我们需要把"无标量"条件改写为可施加于任意矢量函数 $\bA(x,y,z,t)$ 的形式。对于平面波 $\bA(x,y,z,t)=\bA_{0}\sin(\bk\cdot\bx)$（其中要求 $\bk\cdot\bA_{0}=0$），可见

\[
\partial_{x}A^{x}+\partial_{y}A^{y}+\partial_{z}A^{z}+\partial_{t}A^{t}
=(A_{0}^{x}k^{x}+A_{0}^{y}k^{y}+A_{0}^{z}k^{z}+A_{0}^{t}k^{t})\cos(\bk\cdot\bx)
=(\bk\cdot\bA_{0})\cos(\bk\cdot\bx)=0
\]

所以我们宣布，下面一对方程构成\textbf{黎曼矢量波方程}（Riemannian Vector Wave equations, RVW）：

\begin{align}
\partial_{x}^{2}\bA+\partial_{y}^{2}\bA+\partial_{z}^{2}\bA+\partial_{t}^{2}\bA+\wms^{2}\bA &= \bzero \tag{RVW}\\[2pt]
\partial_{x}A^{x}+\partial_{y}A^{y}+\partial_{z}A^{z}+\partial_{t}A^{t} &= 0 \tag{Transverse}
\end{align}

其中第一式告诉我们 $\bA$ 的四个分量各自满足 RSW 方程，而\textbf{横波条件}（transverse condition）则保证这个波是纯矢量波。

对于不介意数学水平略高一些的读者，此主题还有进一步的补充材料。

\section{对称群的表示（Representations of the Symmetry Groups）}

在关于对称性的注记中，我们识别了欧几里得四维空间的若干对称群。其中最大的群是四维欧几里得群 $E(4)$，它包含一切保持点间距离的变换。该群之内有若干子群，包括：

\begin{itemize}
  \item 一切\emph{平移}组成的群，记作 $T(4)$；
  \item 正交群 $O(4)$，即一切保持原点不动的\emph{旋转与反射}；
  \item 特殊正交群 $SO(4)$，即一切保持原点不动的\emph{旋转}。
\end{itemize}

这些群中的变换可以从两个角度来看。若欧几里得四维空间中进行的各种过程服从物理定律——无论谈论的是经典粒子碰撞、波从一处传播到另一处，还是某个量子力学过程——那么我们要求的定律的对称性应当保证：把 $E(4)$ 中的任何变换施加于原来的过程，得到的新过程也服从定律。举一个简单的例子：若设想若干个粒子的世界线以服从能量-动量守恒的方式碰撞，那么我们应当能够把这些世界线一起旋转、反射和平移，而得到一个同样服从守恒律的新碰撞。

另一种看待这些对称性的方式，是把它们视为从一个坐标系到另一个坐标系的变换。若我使用某个坐标系、你使用另一个，只要我们都使用直角坐标——即 $x,y,z,t$ 四根轴彼此正交、四根轴采用同一尺度——那么我把 $\mathbb{R}^{4}$ 中的坐标四元组指定给四维空间中的实际物理点的方式会与你的不同，但我们两套坐标之间由一个 $E(4)$ 中的对称变换相联系。若我们描述同样的物理过程、只是给事物指定了不同的坐标，那么我们都应能对观测施加物理定律，并在任何客观问题上得出相同结论。某个碰撞中能量-动量是否守恒，这是个客观事实——我们两人用各自的坐标系描述同一组事件，应当对该问题得出相同结论。

无论从哪个角度，理解欧几里得空间的对称性都有助于我们理解能满足这些对称性的定律类型。而在理解我们自身宇宙的物理时，最有力的工具之一被证明是数学中称为\textbf{表示论}（representation theory）的分支。

\subsection{表示（Representations）}

设有群 $G$ 与矢量空间 $V$。考虑从 $V$ 到 $V$ 的线性函数（常称为"线性算符"），取其中一切有逆的函数，以函数复合为群乘法，可构成一个群，称为 $V$ 上的一般线性群 $GL(V)$。$G$ 在 $V$ 上的一个\textbf{表示}（representation）是一个函数 $\rho:G\to GL(V)$，满足对任意 $g,h\in G$：

\[
\rho(g)\rho(h)=\rho(gh)
\]

换言之，以下两个过程

\begin{itemize}
  \item 分别对 $g$ 和 $h$ 施加 $\rho$，再用 $GL(V)$ 中的群乘法合并结果；
  \item 先在 $G$ 中把 $g$ 与 $h$ 相乘，再用函数 $\rho$ 把 $gh$ 对应到 $GL(V)$ 的元素；
\end{itemize}

必须给出相同的最终结果，$\rho$ 才是 $G$ 的一个\emph{表示}。换一种说法，$\rho$ 是从 $G$ 到 $GL(V)$ 的一个群同态。

有些群本身就被定义（或等价于）$n\times n$ 矩阵群，例如 $SO(3)$ 与 $SO(4)$。此时它们显然在 $\mathbb{R}^{n}$ 与 $\mathbb{C}^{n}$ 上有表示（若矩阵为实值），或仅在 $\mathbb{C}^{n}$ 上有表示（若矩阵为复值），其中矢量空间上的线性算符就是矩阵与矢量的乘法。这称为该群的\textbf{基本表示}（fundamental representation）。

下面是另一种表示的实例。实数集合 $\mathbb{R}$ 以\emph{加法}为群运算构成一个群，记作 $(\mathbb{R},+)$。取二维矢量空间 $\mathbb{R}^{2}$，则 $GL(\mathbb{R}^{2})$ 由一切从 $\mathbb{R}^{2}$ 到 $\mathbb{R}^{2}$ 的可逆线性函数组成，也可以把它们看成有逆的 $2\times2$ 矩阵。定义表示 $\rho_{1}$ 如下：对 $(\mathbb{R},+)$ 中任意数 $r$，令 $\rho_{1}(r)$ 为把 $\mathbb{R}^{2}$ 中矢量旋转角度 $r$ 的矩阵：

\begin{equation}
\rho_{1}(r)=\begin{pmatrix}\cos r & -\sin r\\ \sin r & \cos r\end{pmatrix} \tag{1}
\end{equation}

不难看出 $\rho_{1}(r+s)=\rho_{1}(r)\rho_{1}(s)$，其中左边用的是 $(\mathbb{R},+)$ 中的加法，右边用的是 $GL(\mathbb{R}^{2})$ 中的矩阵乘法。其实无须真正做矩阵乘法就知道：先旋转角度 $s$ 再旋转角度 $r$ 等于旋转角度 $r+s$。

\subsection{可约与不可约表示}

在上面给出的 $(\mathbb{R},+)$ 在 $\mathbb{R}^{2}$ 上的表示 $\rho_{1}$（由式 (1) 定义）中，很明显我们"用满"了整个矢量空间：不存在 $\mathbb{R}^{2}$ 的某个子空间，使得当 $\rho_{1}$ 对群 $(\mathbb{R},+)$ 各元素给出的矩阵作用其上时，其中的矢量仍留在其中。

但现在考虑 $(\mathbb{R},+)$ 在 $\mathbb{R}^{2}$ 上的另一个表示：

\begin{equation}
\rho_{2}(r)=\begin{pmatrix} \exp(r) & 0\\ 0 & \exp(2r)\end{pmatrix} \tag{2}
\end{equation}

当乘以这样的矩阵时，$\mathbb{R}^{2}$ 中位于 $x$ 轴上的矢量仍留在 $x$ 轴上，位于 $y$ 轴上的矢量仍留在 $y$ 轴上。我们把表示 $\rho_{2}$ 称为\textbf{可约}（reducible），因为它实际上是由两个各自一维的表示"粘在一起"，而不是一个真正二维的表示。

与此相对，我们把先前的表示 $\rho_{1}$ 称为\textbf{不可约}（irreducible），因为不存在 $\mathbb{R}^{2}$ 中被所有 $\rho_{1}(r)$ 的作用保持的子空间。〔值得指出，同样的矩阵作用在\emph{复}数对 $\mathbb{C}^{2}$ 上时给出的是一个可约表示；用实数还是复数可以造成根本差别。〕

不可约表示是构造其他表示的基本积木，因此从纯数学角度值得关注。而在物理学中，事实证明每一个\emph{基本粒子}都对应着相关对称群的一个不同的\emph{不可约}表示！在关于矢量波的注记中，当我们添加一个条件以确保正确区分矢量波与标量波时，所做的正是确保我们挑出了一个不可约表示。

\subsection{酉表示（Unitary Representations）}

量子力学用复矢量空间中的矢量描述物理系统的态。量子态的矢量空间总是配备所谓的\textbf{内积}（inner product），它是点积概念向复数的推广。例如，在复数对的矢量空间 $\mathbb{C}^{2}$ 上，标准内积为

\[
\langle(a,b),(c,d)\rangle=a^{*}c+b^{*}d
\]

这与 $\mathbb{R}^{2}$ 上的标准点积相同，只是对第一个矢量的分量取复共轭。（有时内积按对第二个矢量取复共轭来定义，而非第一个。究竟共轭哪一个是约定问题，真正要紧的只是保持一致。）

注意这会给我们一个非负的\emph{实数}，尽管 $\bv$ 的分量可以是任意复数。我们这样使用内积，以给每个矢量指定长度或模 $|\bv|$。例如

\[
|(a,b)|^{2}=\langle(a,b),(a,b)\rangle=a^{*}a+b^{*}b=|a|^{2}+|b|^{2}
\]

其中 $|a|,|b|$ 指复数 $a,b$ 的绝对值。

若一个带内积的复矢量空间同时也是完备的，则称为希尔伯特空间（Hilbert space）。量子力学中使用的复矢量空间总是希尔伯特空间。

若两个单位矢量 $\bv$ 与 $\bw$ 表示某物理系统的两个不同量子态，则系统被制备在态 $\bv$ 时、通过"处于态 $\bw$"的检验的概率为

\[
P(\text{通过 }\bw\text{ 检验}\mid\text{制备为 }\bv)=|\langle\bv,\bw\rangle|^{2}
\]

现在，正如我们称保持点积的矩阵为\emph{正交}矩阵，我们称保持复内积的矩阵或线性算符为\textbf{酉}（unitary）的：

\[
\langle U(\bv),U(\bw)\rangle=\langle\bv,\bw\rangle
\]

如何判断一个矩阵描述的是不是酉算符？对于正交矩阵，条件是 $M^{T}=M^{-1}$，即矩阵的转置等于其逆。矩阵为酉的条件非常相似，只是用厄米伴随代替转置，即转置的复共轭，记作 $M^{*}$，矩阵为酉当且仅当 $M^{*}=M^{-1}$。

群 $G$ 在希尔伯特空间 $V$ 上的表示 $\rho$ 称为酉表示，若对每个 $g\in G$，$\rho(g)$ 都是 $V$ 上的酉线性算符。

对称群的酉表示在量子力学中扮演特殊角色，因为它们在保持内积的同时也保持概率。若你我在四维空间上使用不同的坐标，那么对于某个量子系统的同一组物理态，我们一般会用系统希尔伯特空间 $V$ 中\emph{不同}的复矢量来描述。但 $V$ 上应当存在相应对称群的一个表示，使我把你的矢量换算成我的——若该表示是\emph{酉}的，我们就会对矢量间所有内积、从而对所预言的概率达成一致。

不过实际上，由于概率都形如 $|\cdot|^{2}$，即便我的内积与你的内积相差一个绝对值为 1 的因子，二者仍然一致。绝对值为 1 的复数 $z$（$|z|=1$）称为一个\textbf{相位}（phase），所以我们也可以说我们的内积"相差一个相位"。任何纯虚数的指数其绝对值为 1，故相位常写成 $\exp(\ii\theta)$，$\theta$ 为实数。满足"通常表示定义的右端多乘一个相位"这一关系的函数 $\rho_{P}:G\to GL(V)$：

\[
\rho_{P}(g)\rho_{P}(h)=\exp\bigl(\ii\theta(g,h)\bigr)\,\rho_{P}(gh)
\]

称为 $G$ 在 $V$ 上的\textbf{投影表示}（projective representation）。

在我们自己的宇宙中确实存在一些物理系统服从相关对称群的投影表示而非"真正的"表示，我们也预期黎曼宇宙中会发生同样的事。不过这里不打算对处理投影表示的微妙之处给出严格论述，因为我们主要通过处理\textbf{双覆盖}（double cover）——即出现此问题的两个群的二重覆盖——来大体绕开它们。

\subsection{SO(3) 与 SO(4) 的双覆盖（Double Covers）}

在我们宇宙中，若你把坐标系绕某根轴连续旋转角度 $\theta$，则你用来描述自旋与该轴对齐的电子的态矢量将被乘以相位 $\exp(\ii\theta/2)$。当 $\theta$ 达到 $2\pi$、坐标系回到原处时，电子的态矢量已改变了一个因子 $\exp(\ii\pi)=-1$。这是因为电子的态矢量服从三维旋转群 $SO(3)$ 的\emph{投影表示}，而非真正的表示。

\subsubsection{SO(3) 与 SO(4) 的拓扑}

\orthofig{images/03__009.png}{SO(3) 中的不可收缩回路}
\orthofig{images/03__010.png}{SO(3) 中的可收缩回路}

描绘 $SO(3)$ 的一种方式，是把它想象成一个半径为 $\pi$ 的实心三维球，球心为 $O$。球内任一点 $P$ 可表示一个旋转，旋转角等于距离 $OP$，转轴为线段 $OP$；旋转方向取为：右手拇指从 $O$ 指向 $P$ 时四指卷曲的方向。在此方案下，球面上相对的两点表示相同的旋转，因为绕同一轴朝两个方向旋转 $\pi$ 效果相同。于是该球的任一直径实际上都是 $SO(3)$ 中的一条\emph{回路}：它的两个端点在球面上对应完全相同的旋转，故应当视为同一点，从而使回路闭合。但与普通空间中的回路不同，这条回路无法连续收缩为一点。直径的两个端点可以一起移动，只要始终保持彼此相对，但它们不能靠拢。

$SO(3)$ 中的回路分两类：可以连续形变收缩为一点的\textbf{可收缩}（contractible）回路，以及不能收缩的回路。把绕某轴的旋转角从 0 连续增大到 $2\pi$，得到一条不可收缩回路，它在我们的旋转球图像中就是一条直径——例如\textbf{右上图中的蓝色直线}。然而若继续把旋转角增大到 $4\pi$，就得到一条可收缩回路！\textbf{右下图中}近乎笔直的那对绿线说明了这一点（为清晰起见它们避免重复任何旋转）；沿绿—黄—红的顺序看，可见该路径收缩为一点。

我们说 $SO(3)$ 不是\textbf{单连通}（simply connected）的。在单连通拓扑空间中，从点 $A$ 到点 $B$ 的任意两条连续路径都可以连续地形变为彼此。在 $SO(3)$ 中，从 $A$ 到 $B$ 的路径总有两类，其中一类的路径无法形变为另一类的路径。

给定一个群的投影表示，我们总可以尝试通过把 $\rho_{P}(g)$ 再乘一个相位，把它变成真正的表示；这样做不会改变由矢量间内积得到的概率。但事实证明，只有当群是\emph{单连通}的，我们才能保证相位可以被完全消去。对于 $SO(3)$ 这样的群，我们不得不与某些投影表示打交道。

处理此问题的一个数学上方便的招数，是构造一个新群 $SO(3)^{\sim}$，它由从 $SO(3)$ 中的恒等元到群中各元素的\emph{路径等价类}组成，而不仅仅是群元素本身。若 $p$ 是某条具体路径，用 $[p]$ 表示一切与 $p$ 有相同端点、且可连续形变为 $p$ 的路径所构成的等价类。

若取从 $I_{3}$ 到 $SO(3)$ 中 $g$ 的一切连续路径，然后把可连续形变为彼此的路径视为等价，我们恰好得到从 $I_{3}$ 到 $g$ 的两个路径等价类。所以对 $SO(3)$ 中每个 $g$，新群 $SO(3)^{\sim}$ 都含有两个元素 $g_{1},g_{2}$，它们可视为从 $SO(3)$ 的恒等元连续移动到 $g$ 的两种不同方式。

如何在 $SO(3)^{\sim}$ 中相乘两个元素？形式上，一条穿过 $SO(3)$ 的路径是函数 $p:[0,1]\to SO(3)$，其中 $[0,1]$ 是 0 到 1 之间（含端点）的实数集合。取所要相乘的两个等价类的代表路径 $p_{1}$ 与 $p_{2}$，定义 $[p_{1}]$ 与 $[p_{2}]$ 之积为路径 $p_{3}(t)=p_{1}(t)\,p_{2}(t)$ 的等价类。这同时刻画了旋转的乘法与"路径类的乘法"的细节。

在 $SO(3)$ 中乘以路径类比听起来容易：对固定端点而言，它恰好像二元群 $\{1,-1\}$ 中的乘法，其中任何元素自乘都给出恒等元。

一个具体例子：考虑绕 $z$ 轴旋转角度从 0 到 $2\pi$ 所得到的路径

\[
p(t)=\begin{pmatrix}\cos(2\pi t) & -\sin(2\pi t) & 0\\ \sin(2\pi t) & \cos(2\pi t) & 0\\ 0 & 0 & 1\end{pmatrix}
\]

它是一条从恒等元回到恒等元的回路（即 $p(0)=p(1)=I_{3}$），但它\emph{不能}连续形变为恒等路径

\[
e(t)=\begin{pmatrix}1 & 0 & 0\\ 0 & 1 & 0\\ 0 & 0 & 1\end{pmatrix}
\]

然而 $[p][p]$ 是等价类 $[p^{2}]$，其中

\[
p^{2}(t)=\begin{pmatrix}\cos(4\pi t) & -\sin(4\pi t) & 0\\ \sin(4\pi t) & \cos(4\pi t) & 0\\ 0 & 0 & 1\end{pmatrix}
\]

该路径绕 $z$ 轴旋转角度从 0 到 $4\pi$，而它\emph{能}连续形变为恒等路径。所以 $[p][p]=[p^{2}]=[e]$。

不过，与其完全按上述定义处理 $SO(3)^{\sim}$，不如找一个与之同构的矩阵群来得更简单。暂时把这个尚未识别的矩阵群称为"$G_{3}$"——我们把 $G_{3}$ 与 $SO(3)^{\sim}$ 本身都称为 $SO(3)$ 的\textbf{双覆盖}（double cover）。双覆盖将是单连通群，所以我们只需处理它的真表示；但这样做时，我们实际上同时处理了 $SO(3)$ 的真表示\emph{与}投影表示。

为说明 $G_{3}$ 是 $SO(3)$ 的双覆盖，我们需要一个 2 对 1 的函数 $f_{3}:G_{3}\to SO(3)$，使得对一切 $g_{1},g_{2}\in G_{3}$：

\[
f_{3}(g_{1})f_{3}(g_{2})=f_{3}(g_{1}g_{2})
\]

这表明 $G_{3}$ 中的乘法刻画了 $SO(3)$ 中旋转乘法的含义，是群同态的一个例子。为表明 $G_{3}$ 中的乘法同时刻画了路径类乘法这一（非常简单的）结构，我们要求 $G_{3}$ 中被 $f_{3}$ 映到 $SO(3)$ 中恒等元 $I_{3}$ 的那两个元素：

\[
f_{3}(i_{1})=f_{3}(i_{2})=I_{3}
\]

的乘法表恰好与二元群 $\{1,-1\}$ 相同。换句话说，其中一个是 $G_{3}$ 中的恒等元，另一个自乘则给出恒等元。

注意，对 $SO(3)^{\sim}$ 本身我们很容易构造一个类似的函数 $f_{3^{\sim}}:SO(3)^{\sim}\to SO(3)$：对任意路径 $q:[0,1]\to SO(3)$，定义 $f_{3^{\sim}}([q])=q(1)$，而 $f_{3^{\sim}}$ 映到 $SO(3)$ 中恒等元的两个元素就是上面定义的 $[p]$ 与 $[e]$。

群 $SO(4)$ 也不是单连通的，且具有完全相同的诸多问题，因此我们也想为 $SO(4)$ 找一个双覆盖 $G_{4}$ 以及类似的同态 $f_{4}:G_{4}\to SO(4)$。$SO(4)$ 的拓扑是 $SO(3)\times S^{3}$：即它与 $SO(3)$ 与三维球面的直积同胚。为看出这点，对 $SO(4)$ 中给定元素 $R$，考虑它把 $t$ 方向单位矢量映到的矢量 $R\be_{t}$。一切这样的矢量将构成 $\mathbb{R}^{4}$ 中的单位三维球面，因为显然存在某个旋转把 $\be_{t}$ 映到 $\mathbb{R}^{4}$ 中任意单位矢量。若固定某个单位矢量 $\bu$，则 $SO(4)$ 中把 $\be_{t}$ 映到 $\bu$ 的那些元素 $R$ 构成 $SO(4)$ 的一个子群；不同 $\bu$ 对应的各子群彼此同构，而对 $\bu=\be_{t}$ 的选择，该子群按定义就是 $SO(3)$。这不是关于 $SO(4)$ 拓扑断言的严格证明，但应能使它直观可信。

\subsubsection{构造 2 对 1 同态}

为识别 $SO(3)$ 与 $SO(4)$ 的双覆盖，我们从四个 $2\times2$ 矩阵出发：

\[
H_{x}=\begin{pmatrix}0 & -\ii\\ -\ii & 0\end{pmatrix},\quad
H_{y}=\begin{pmatrix}0 & -1\\ 1 & 0\end{pmatrix},\quad
H_{z}=\begin{pmatrix}-\ii & 0\\ 0 & \ii\end{pmatrix},\quad
H_{t}=\begin{pmatrix}1 & 0\\ 0 & 1\end{pmatrix}
\]

这四个矩阵的行列式都为 1。把其中任意两个相乘，得到的要么还是这个集合中的某个矩阵，要么是其相反数。乘法表如下，其中取最左列的矩阵右乘顶行的矩阵。〔若这个乘法表看起来眼熟，也许是因为它给出了四元数（quaternions）的一个矩阵表示。〕

\begin{center}
\begin{tabular}{c|cccc}
 & $H_{x}$ & $H_{y}$ & $H_{z}$ & $H_{t}$\\\hline
$H_{x}$ & $-H_{t}$ & $H_{z}$ & $-H_{y}$ & $H_{x}$\\
$H_{y}$ & $-H_{z}$ & $-H_{t}$ & $H_{x}$ & $H_{y}$\\
$H_{z}$ & $H_{y}$ & $-H_{x}$ & $-H_{t}$ & $H_{z}$\\
$H_{t}$ & $H_{x}$ & $H_{y}$ & $H_{z}$ & $H_{t}$
\end{tabular}
\end{center}

对 $\mathbb{R}^{4}$ 中任一点 $\bx=(x,y,z,t)$，定义 $H(\bx)$ 为

\[
H(\bx)=x^{\mu}H_{\mu}=xH_{x}+yH_{y}+zH_{z}+tH_{t}
=\begin{pmatrix} t-z\ii & -y-x\ii\\ y-x\ii & t+z\ii\end{pmatrix}
\]

我们可以把 $2\times2$ 复矩阵的集合看成 8 维\emph{实}矢量空间：显然可以相加这样的矩阵、用实数相乘，其中一个基是八个在矩阵四个位置之一处为 1 或 $\ii$、其余为零的矩阵。于是 $H$ 是 $\mathbb{R}^{4}$ 与该 8 维矢量空间中一个四实维子空间之间的同构；我们把这个子空间也称为 $H$。由 $\{H_{x},H_{y},H_{z},H_{t}\}$ 的乘法表可见，$H$ 在矩阵乘法下封闭：$H$ 中任意两元素的矩阵乘积仍在 $H$ 中。

$H(\bx)$ 的行列式是 $\bx$ 的长度的平方：

\[
\det H(\bx)=(t-z\ii)(t+z\ii)-(-y-x\ii)(y-x\ii)=x^{2}+y^{2}+z^{2}+t^{2}=|\bx|^{2}
\]

而 $H(\bx)$ 的迹（对角线元素之和）是 $\bx$ 的 $t$ 坐标的两倍：

\[
\tr H(\bx)=(t-z\ii)+(t+z\ii)=2t
\]

$H(\bx)$ 的厄米伴随也在 $H$ 中。具体地，$H(\bx)^{*}=H((-x,-y,-z,t))$：

\[
H(\bx)^{*}=\begin{pmatrix} t+z\ii & y+x\ii\\ -y+x\ii & t-z\ii\end{pmatrix}
\]

而 $H(\bx)$ 与其厄米伴随之积是 $\bx$ 长度平方乘以 $2\times2$ 单位矩阵：

\[
H(\bx)H(\bx)^{*}=\begin{pmatrix}x^{2}+y^{2}+z^{2}+t^{2} & 0\\ 0 & x^{2}+y^{2}+z^{2}+t^{2}\end{pmatrix}=|\bx|^{2}I_{2}
\]

这意味着当 $|\bx|=1$ 时，$H(\bx)$ 是\emph{酉}的。事实上，一切行列式为 1 的 $2\times2$ 酉矩阵构成的群，即 \textbf{SU(2)}，正是 $H$ 中由 $\mathbb{R}^{4}$ 里单位矢量 $\bu$ 对应的矩阵 $H(\bu)$ 构成的子集。

除了 $\bx$ 的模，我们还可以从 $\bx,\by$ 在 $H$ 中的像提取它们的点积：

\[
\bx\cdot\by=\tfrac{1}{2}\tr\!\bigl(H(\bx)H(\by)^{*}\bigr)
\]

此结果可通过逐项验证 $H$ 矩阵（在上指标表示厄米伴随）得到：

\[
\tfrac{1}{2}\tr\!\bigl(H_{\mu}H^{\nu*}\bigr)=\delta_{\mu}^{\nu}
\]

这也使我们能把映射 $H$ 的逆写为

\[
H^{-1}(h)=\tfrac{1}{2}\tr\!\bigl(hH^{\nu*}\bigr)\,\be_{\nu}
\]

其中 $\{\be_{\nu}\}$ 是 $\mathbb{R}^{4}$ 的坐标基。

设 $g,h\in SU(2)$。我们如下定义函数 $f_{4}:SU(2)\times SU(2)\to SO(4)$——即从 $SU(2)$ 元素\emph{对}的集合到 $SO(4)$ 的函数：

\[
f_{4}(g,h)\,\bx=H^{-1}\!\bigl(g\,H(\bx)\,h^{-1}\bigr)
\]

由于 $g$ 与 $h^{-1}$ 都在 $SU(2)\subseteq H$ 中，乘积 $gH(\bx)h^{-1}$ 也在 $H$ 中，所以 $f_{4}(g,h)$ 是 $\mathbb{R}^{4}$ 上良定义的线性算符。利用 $\det(H(\bx))$ 与 $|\bx|^{2}$ 的关系，以及"矩阵乘积的行列式等于行列式之积"：

\[
|f_{4}(g,h)\bx|^{2}=\det\!\bigl(H(f_{4}(g,h)\bx)\bigr)
=\det\!\bigl(gH(\bx)h^{-1}\bigr)
=\det(g)\det\!\bigl(H(\bx)\bigr)\det(h^{-1})
=\det\!\bigl(H(\bx)\bigr)=|\bx|^{2}
\]

所以 $f_{4}(g,h)$ 把 $\mathbb{R}^{4}$ 中矢量映为长度相同的矢量，从而它要么是旋转、要么是反射。而结果总是旋转，即 $SO(4)$ 的元素。

$f_{4}$ 能给出 $SO(4)$ 的\emph{每一个}元素吗？稍作计算便不难表明，$f_{4}(g,h)$ 能给出六个坐标面内任意角度的旋转。通过复合这种形式的旋转，可以构造出任意旋转。

\begin{center}
\begin{tabular}{lll}
\toprule
$g$ & $h$ & 旋转所在的坐标面（角 $2\theta$）\\
\midrule
$\cos\theta\,H_{t}+\sin\theta\,H_{x}$ & $\cos\theta\,H_{t}+\sin\theta\,H_{x}$ & $yz$ 面\\
$\cos\theta\,H_{t}+\sin\theta\,H_{y}$ & $\cos\theta\,H_{t}+\sin\theta\,H_{y}$ & $xz$ 面\\
$\cos\theta\,H_{t}+\sin\theta\,H_{z}$ & $\cos\theta\,H_{t}+\sin\theta\,H_{z}$ & $xy$ 面\\
$\cos\theta\,H_{t}-\sin\theta\,H_{x}$ & $\cos\theta\,H_{t}+\sin\theta\,H_{x}$ & $xt$ 面\\
$\cos\theta\,H_{t}-\sin\theta\,H_{y}$ & $\cos\theta\,H_{t}+\sin\theta\,H_{y}$ & $yt$ 面\\
$\cos\theta\,H_{t}-\sin\theta\,H_{z}$ & $\cos\theta\,H_{t}+\sin\theta\,H_{z}$ & $zt$ 面\\
\bottomrule
\end{tabular}
\end{center}

我们按显然的方式把 $SU(2)\times SU(2)$ 变成一个群，即把对中各元素分别相乘：$(g_{1},h_{1})(g_{2},h_{2})=(g_{1}g_{2},h_{1}h_{2})$。用此乘法可得

\begin{align*}
f_{4}\bigl((g_{1},h_{1})(g_{2},h_{2})\bigr)\bx
&=f_{4}(g_{1}g_{2},h_{1}h_{2})\bx\\
&=H^{-1}\!\bigl(g_{1}g_{2}H(\bx)h_{2}^{-1}h_{1}^{-1}\bigr)\\
&=H^{-1}\!\bigl(g_{1}H\bigl(H^{-1}(g_{2}H(\bx)h_{2}^{-1})\bigr)h_{1}^{-1}\bigr)\\
&=f_{4}(g_{1},h_{1})f_{4}(g_{2},h_{2})\bx
\end{align*}

同时还有

\[
f_{4}(I_{2},I_{2})=I_{4},\qquad f_{4}(-I_{2},-I_{2})=I_{4}
\]

这表明 $SU(2)\times SU(2)$ 中有两个元素映到 $SO(4)$ 的恒等元，而这两个元素具有与 $\{1,-1\}$ 相同的群结构。这就满足了为表明 $SU(2)\times SU(2)$ 是 $SO(4)$ 的双覆盖所需要的条件。

我们是从讨论 $SO(3)$（而非 $SO(4)$）的双覆盖开始的，但上述一切都能轻易改造成此目的。若如下定义 $f_{3}:SU(2)\to SO(3)$：

\[
f_{3}(g)\,\bx=f_{4}(g,g)\,\bx=H^{-1}\!\bigl(g\,H(\bx)\,g^{-1}\bigr)
\]

则 $f_{3}(g)$ 仍在 $\mathbb{R}^{4}$ 上表现为旋转。但假设我们把 $\bx$ 限制到 $\mathbb{R}^{3}$，即令 $t$ 坐标为 0。由于 $t$ 与 $H(\bx)$ 的迹之间的关系，这意味着 $\tr H(\bx)=0$，再由迹的循环性质可得

\[
\tr H(f_{3}(g)\bx)=\tr\!\bigl(gH(\bx)g^{-1}\bigr)=\tr\!\bigl(g^{-1}gH(\bx)\bigr)=\tr H(\bx)=0
\]

所以 $f_{3}(g)\bx$ 的 $t$ 坐标也为 0。既然一切从 $\mathbb{R}^{4}$ 中 $t=0$ 三维子空间出发的矢量都留在该子空间内，$f_{3}(g)$ 给出三维中的旋转，即 $SO(3)$ 的元素。

总结起来，我们已确立：

\begin{itemize}
  \item $SU(2)$ 是 $SO(3)$ 的双覆盖；
  \item $SU(2)\times SU(2)$ 是 $SO(4)$ 的双覆盖。
\end{itemize}

\subsection{SU(2) 的表示}

既然已看到 $SU(2)$ 与 $SO(3)$、$SO(4)$ 的联系，下一步就是构造 $SU(2)$ 的不可约酉表示。由这些表示，我们可以进一步构造 $SO(3)$ 与 $SO(4)$ 的真表示和投影表示。

\subsubsection{自旋-0}

$SU(2)$ 最简单的不可约酉表示是每个群都有的一维\textbf{平凡表示}（trivial representation）：

\[
\rho_{0}(g)=I_{1}
\]

复数 $\mathbb{C}$ 是一个一复维矢量空间，我们与 $SU(2)$ 中每个元素 $g$ 关联的线性算符是 $1\times1$ 单位矩阵。这称为\textbf{自旋-0}（spin-0）表示，它描述标量在坐标变换下的变换方式：完全不变！

\subsubsection{自旋-$\tfrac12$}

$SU(2)$ 次简单的不可约酉表示是二维的：

\[
\rho_{\half}(g)=g
\]

换言之，我们直接用 $SU(2)$ 元素 $g$，通过通常的矩阵与矢量乘法，在 $\mathbb{C}^{2}$ 上定义线性算符。这称为\textbf{自旋-$\half$}（spin-$\half$）表示。

\subsubsection{自旋-1}

我们可以如下在 $\mathbb{C}^{3}$ 上构造 $SU(2)$ 的一个三维不可约酉表示：

\[
\rho_{1}(g)=f_{3}(g)
\]

我们用来表明 $SU(2)$ 是 $SO(3)$ 双覆盖的函数 $f_{3}$ 给出一个 $SO(3)$ 的元素，但没有理由不能把这个 $3\times3$ 实值矩阵当作复数三元组空间 $\mathbb{C}^{3}$ 上的线性算符。这称为\textbf{自旋-1}（spin-1）表示。

\subsubsection{自旋-$j$}

接下来，我们对任意半整数 $j$ 给出\textbf{自旋-$j$}（spin-$j$）表示的一般构造；这是 $SU(2)$ 的一个 $(2j+1)$ 维不可约酉表示。

从任意两个维数为 $n$ 与 $m$ 的矢量空间 $V$ 与 $W$，可以构造一个新矢量空间：$n\times m$ 矩阵的空间。矩阵可按显然方式相加、用实数或复数相乘。这个新矢量空间称为 $V$ 与 $W$ 的张量积（tensor product），记作 $V\otimes W$。

若 $V$ 有基 $\{\be_{1},\dots,\be_{n}\}$、$W$ 有基 $\{\be'_{1},\dots,\be'_{m}\}$，我们定义 $\be_{i}\otimes\be'_{j}$ 为一个 $n\times m$ 矩阵，除第 $i$ 行第 $j$ 列为 1 外其余皆为零。所有 $\be_{i}\otimes\be'_{j}$ 构成 $V\otimes W$ 的基。再定义 $V$ 中任意矢量 $\bv=v^{i}\be_{i}$ 与 $W$ 中任意矢量 $\bw=w^{j}\be'_{j}$ 之间的张量积为

\[
\bv\otimes\bw=v^{i}w^{j}\,\be_{i}\otimes\be'_{j}
\]

其中重复指标采用爱因斯坦求和约定。也就是说，我们把 $\bv$ 的分量竖排在一个 $n\times m$ 矩阵左侧、把 $\bw$ 的分量横排在该矩阵上方，再把矩阵中位于同一行同一列的矢量分量之积填入。当然 $V\otimes W$ 的一般元素并非都取这种形式，但总能把它写成张量积的线性组合。

可以如法构造任意多个矢量空间的张量积 $U\otimes V\otimes W\otimes\cdots$，得到可看作向三维以上延展的类矩阵数组的空间。若各矢量空间上都定义了内积，则可通过把各个内积相乘来定义张量积上的内积：

\[
\langle\bu\otimes\bv,\bw\otimes\bx\rangle=\langle\bu,\bw\rangle\langle\bv,\bx\rangle
\]

此内积通过在第一宗量上线性、在第二宗量上共轭线性，延拓到\emph{一切}张量对。

现在，我们可以在 $\mathbb{C}^{2}$ 的 $2j$ 重张量积上构造 $SU(2)$ 的一个表示：

\[
\rho_{\half\otimes 2j}(g)\,v_{1}\otimes v_{2}\otimes\cdots\otimes v_{2j}
=\rho_{\half}(g)v_{1}\otimes\rho_{\half}(g)v_{2}\otimes\cdots\otimes\rho_{\half}(g)v_{2j}
=gv_{1}\otimes gv_{2}\otimes\cdots\otimes gv_{2j}
\]

这就定义了 $\rho_{\half\otimes 2j}(g)$ 在 $\mathbb{C}^{2}$ 的 $2j$ 个元素的张量积上的作用，而 $\mathbb{C}^{2}\otimes\mathbb{C}^{2}\otimes\cdots$ 中其他元素的结果由线性性给出。注意

\[
\rho_{\half\otimes 2j}(-g)=(-1)^{2j}\rho_{\half\otimes 2j}(g)
\]

这是 $SU(2)$ 的一个酉表示，但当 $j\ge 1$ 时它\emph{不是}不可约的。例如 $j=1$ 时，$\mathbb{C}^{2}\otimes\mathbb{C}^{2}$ 中由反对称 $2\times2$ 矩阵构成的一维子空间在任何 $g$ 的 $\rho_{\half\otimes\half}(g)$ 作用下保持不变。设反对称矩阵为

\[
\begin{pmatrix}0 & a\\ -a & 0\end{pmatrix}=a\bigl(\be_{1}\otimes\be_{2}-\be_{2}\otimes\be_{1}\bigr)
\]

则

\[
\rho_{\half\otimes\half}(g)\bigl[a(\be_{1}\otimes\be_{2}-\be_{2}\otimes\be_{1})\bigr]
=a\,(g\be_{1})\otimes(g\be_{2})-a\,(g\be_{2})\otimes(g\be_{1})
\]

把 $g=\begin{pmatrix}p & q\\ r & s\end{pmatrix}$ 代入（$g\be_{1}=(p,r)$，$g\be_{2}=(q,s)$），并以 $\det g=ps-rq=1$ 化简，结果仍为 $\begin{pmatrix}0 & a\\ -a & 0\end{pmatrix}$ 的形式，故该子空间被保持。

所以 $\mathbb{C}^{2}$ 上自旋-$\half$ 表示的 $2j$ 重张量积并不能给出不可约表示。但它总是\emph{包含}一个我们感兴趣的特定不可约表示：作用在\textbf{对称}张量子空间上的那一个。正如我们说矩阵对称是指交换任意分量的指标得到同值的另一分量（$M^{ij}=M^{ji}$），我们说张量对称是指交换任意分量的\emph{任意}一对指标得到同值的另一分量，例如 $T^{ijklm}=T^{jiklm}$。但若能交换任意一对指标，也就能对指标作任意排列，故 $n$ 个指标的全部 $n!$ 种排列都必须给出同值的分量。

一个等价的条件是：对称张量可以写成若干项之和，其中每一项是同一组 $n$ 个矢量的张量积在所有排列上的求和。例如 $n=3$ 时，每一项形如

\[
\bu\otimes\bv\otimes\bw+\bv\otimes\bw\otimes\bu+\bw\otimes\bu\otimes\bv
+\bw\otimes\bv\otimes\bu+\bu\otimes\bw\otimes\bv+\bv\otimes\bu\otimes\bw
\]

对应某个矢量三元组 $\bu,\bv,\bw$。这 $n$ 个矢量不必互异；若它们有重复，某些张量积会在和中出现多次。我们可以简单地通过改变 $\be_{2}$ 在矢量组中出现的次数（从 0 到 $n$）来构造 $\mathbb{C}^{2}$ 的 $n$ 重张量积中对称张量的基，共得 $n+1$ 个矢量。例如 $n=3$ 时，$\mathbb{C}^{2}\otimes\mathbb{C}^{2}\otimes\mathbb{C}^{2}$ 中对称张量的一个正交归一基为

\begin{align*}
\Bigl\{&\be_{1}\otimes\be_{1}\otimes\be_{1},\;
\frac{\be_{2}\otimes\be_{1}\otimes\be_{1}+\be_{1}\otimes\be_{2}\otimes\be_{1}+\be_{1}\otimes\be_{1}\otimes\be_{2}}{\sqrt{3}},\\
&\frac{\be_{1}\otimes\be_{2}\otimes\be_{2}+\be_{2}\otimes\be_{1}\otimes\be_{2}+\be_{2}\otimes\be_{2}\otimes\be_{1}}{\sqrt{3}},\;
\be_{2}\otimes\be_{2}\otimes\be_{2}\Bigr\}
\end{align*}

由于自旋-$\half$ 表示的 $2j$ 重张量积 $\rho_{\half\otimes 2j}$ 只是把 $\rho_{\half}(g)$ 施加于张量积中的每个矢量，它作用在这样一项（排列后张量积之和）上的结果仍是同形式的一项。所以 $\rho_{\half\otimes 2j}$ 把对称张量映为对称张量。

我们定义 $SU(2)$ 的\textbf{自旋-$j$ 表示} $\rho_{j}$ 为 $\rho_{\half\otimes 2j}$ \emph{限制到对称张量子空间上}。该子空间维数为 $2j+1$，我们把它记作 $V_{j}$。

$SU(2)$ 的自旋-$j$ 表示是不可约酉表示，正如从其名称可以猜到，它们对应具有 $j$ 个角动量单位的基本粒子或其他量子力学系统。若 $j$ 为整数，该表示称为\textbf{玻色型}（bosonic）；若 $j$ 为半奇数，则称为\textbf{费米型}（fermionic）。

\subsubsection{等价表示}

本节开头给出了自旋-0、自旋-$\half$ 与自旋-1 表示的定义。显然前两个与我们对自旋-$j$ 的一般定义一致，但自旋-1 的定义看起来完全不同。我们说

\[
\rho_{1}(g)=f_{3}(g)
\]

其中 $f_{3}(g)$ 是 $\mathbb{R}^{3}$ 上的线性算符，定义为

\[
f_{3}(g)\bx=H^{-1}\!\bigl(g\,H(\bx)\,g^{-1}\bigr)
\]

我们随后把相关的实 $3\times3$ 矩阵当成 $GL(\mathbb{C}^{3})$ 的元素。但我们对自旋-$j$ 的一般定义在 $j=1$ 时化为

\[
\rho_{1}(g)\,v\otimes w=gv\otimes gw
\]

其中 $\bv,\bw\in\mathbb{C}^{2}$。我们把定义通过线性性延拓到 $\mathbb{C}^{2}\otimes\mathbb{C}^{2}$ 中任意张量，然后限制到对称张量上。这两个定义怎么可能描述同一件事？

它们显然\emph{不是}在描述\emph{完全}同一件事，但我们说两个表示\textbf{等价}（equivalent），若在两个表示所作用的矢量空间之间存在一个可逆线性函数，它与两个表示对易。也就是说，群 $G$ 在矢量空间 $V_{A}$ 上的表示 $\rho_{A}$ \emph{等价于}同一群在矢量空间 $V_{B}$ 上的表示 $\rho_{B}$，若存在可逆线性函数 $T:V_{A}\to V_{B}$，使得对每个 $g\in G$：

\[
T\rho_{A}(g)=\rho_{B}(g)T
\]

另一种写法是：对每个 $g\in G$ 与每个 $\bv\in V_{A}$，

\[
\rho_{A}(g)\bv=T^{-1}\rho_{B}(g)T\bv
\]

无论 $\rho_{A}(g)$ 对 $V_{A}$ 中 $\bv$ 做了什么，若用 $T$ 把 $\bv$ 送到 $V_{B}$、让 $\rho_{B}(g)$ 作用在其像上、再用 $T^{-1}$ 把结果带回 $V_{A}$，最终得到同一矢量。所以 $T$ 是一种翻译器，用来证明 $\rho_{A}(g)$ 与 $\rho_{B}(g)$ 在做同一件事——只是换了名字、换了地方。

按此定义，2 对 1 同态 $f_{3}:SU(2)\to SO(3)$ 与自旋-1 表示 $\rho_{1}$ 等价。

对量子力学角动量有一定经验的读者不妨尝试下面这道练习：

\begin{itemize}
  \item 把 $SU(2)$ 的一般元素 $g$ 写成 $H((x,y,z,t))$ 的形式，并施加条件 $x^{2}+y^{2}+z^{2}+t^{2}=1$。
  \item 计算 $f_{3}$ 为 $g$ 在标准基下给出的 $SO(3)$ 中的矩阵 $U_{\text{dc}}$。
  \item 计算 $\rho_{1}$ 为 $g$ 在正交归一基 $\{\be_{1}\otimes\be_{1},(\be_{1}\otimes\be_{2}+\be_{2}\otimes\be_{1})/\sqrt{2},\be_{2}\otimes\be_{2}\}$ 下给出的 $V_{1}$ 中的矩阵 $U_{1}$。
  \item 求一个矩阵 $T$ 使 $U_{\text{dc}}=T^{-1}U_{1}T$。\textbf{提示}：把 $V_{1}$ 中的基矢量看作对应 $z$ 轴自旋为 $-1,0,1$ 的量子态，把 $\mathbb{R}^{3}$ 中的基矢量看作对应 $x$ 轴、$y$ 轴、$z$ 轴自旋为 0 的量子态。若你会为由前者构造后者的态（对总自旋为 1 的粒子），那就给出了 $T$。
\end{itemize}

\subsubsection{对偶表示}

设群 $G$ 在某个复矢量空间 $V$ 上有表示 $\rho$。用 $V^{*}$ 表示一切从 $V$ 到复数 $\mathbb{C}$ 的线性函数构成的\emph{对偶矢量空间}。我们如下定义 $G$ 在 $V^{*}$ 上的\emph{对偶表示} $\rho^{*}$：

\[
(\rho^{*}(g)f)(v)=f\bigl(\rho(g^{-1})v\bigr),\qquad \forall f\in V^{*},\;v\in V,\;g\in G
\]

这里是通过指定新函数 $\rho^{*}(g)f$ 在任意矢量 $\bv$ 上的取值来定义它。在此定义中我们需要用 $g$ 的\emph{逆}来对 $\bv$ 施加 $\rho$，以便 $\rho^{*}$ 确实构成一个表示：

\begin{align*}
&(\rho^{*}(gh)f)(v)=f\bigl(\rho((gh)^{-1})v\bigr)\\
&\quad=f\bigl(\rho(h^{-1}g^{-1})v\bigr)
=f\bigl(\rho(h^{-1})\rho(g^{-1})v\bigr)\\
&\quad=(\rho^{*}(h)f)\bigl(\rho(g^{-1})v\bigr)
=\bigl(\rho^{*}(g)\rho^{*}(h)f\bigr)(v)
\end{align*}

结果表明，$SU(2)$ 的一切自旋-$j$ 表示都\emph{等价于}各自的对偶表示。

为说明这点，我们从自旋-$\half$ 开始。$n$ 维列维-奇维塔符号（Levi-Civita symbol）写作 $\varepsilon_{ij\cdots k}$（$n$ 个下标），定义为

\[
\varepsilon_{ij\cdots k}=\begin{cases}+1 & \{i,j,\dots,k\}\text{ 是 }\{1,2,\dots,n\}\text{ 的偶排列}\\ -1 & \{i,j,\dots,k\}\text{ 是 }\{1,2,\dots,n\}\text{ 的奇排列}\\ 0 & \text{其他}\end{cases}
\]

我们如下定义同构 $T_{\half}:\mathbb{C}^{2}\to\mathbb{C}^{2*}$：

\[
\bigl(T_{\half}(v)\bigr)(w)=\varepsilon_{ij}v^{i}w^{j}
\]

为确立 $\rho_{\half}$ 等价于 $\rho_{\half}^{*}$，需证明对一切 $g\in SU(2)$ 与 $\bv\in\mathbb{C}^{2}$，$\mathbb{C}^{2*}$ 中以下两个函数相同：

\[
T_{\half}\bigl(\rho_{\half}(g)v\bigr)=\rho_{\half}^{*}(g)T_{\half}(v)
\]

把它们作用于同一个 $\bw\in\mathbb{C}^{2}$，即成为

\[
\varepsilon_{ik}g^{i}_{j}v^{j}w^{k}=\varepsilon_{jm}(g^{-1})^{m}_{k}v^{j}w^{k}
\]

任一 $2\times2$ 矩阵 $A$ 的行列式可写为

\[
\det A=\varepsilon_{iq}A^{i}_{1}A^{q}_{2}
\]

由于 $g\in SU(2)$ 的行列式为 1，故

\[
\varepsilon_{iq}g^{i}_{1}g^{q}_{2}=1
\]

交换这里的下指标只改变符号：

\[
\varepsilon_{iq}g^{i}_{2}g^{q}_{1}=\varepsilon_{qi}g^{i}_{1}g^{q}_{2}=-\varepsilon_{iq}g^{i}_{1}g^{q}_{2}=-1
\]

重复下指标会给出零，因为 $\varepsilon_{iq}$ 反对称而 $g^{i}_{m}g^{q}_{m}$ 对称：

\[
\varepsilon_{iq}g^{i}_{1}g^{q}_{1}=\varepsilon_{iq}g^{i}_{2}g^{q}_{2}=0
\]

综上所述，这表明确维-奇维塔符号 $\varepsilon$ 在 $g$ 变换下不变：

\[
\varepsilon_{iq}g^{i}_{j}g^{q}_{m}=\varepsilon_{jm}
\]

在右端乘以 $(g^{-1})^{m}_{k}$ 得

\[
\varepsilon_{iq}g^{i}_{j}g^{q}_{m}(g^{-1})^{m}_{k}=\varepsilon_{jm}(g^{-1})^{m}_{k}
\;\Longrightarrow\;
\varepsilon_{iq}g^{i}_{j}\delta^{q}_{k}=\varepsilon_{jm}(g^{-1})^{m}_{k}
\;\Longrightarrow\;
\varepsilon_{ik}g^{i}_{j}=\varepsilon_{jm}(g^{-1})^{m}_{k}
\]

把它作用于 $\mathbb{C}^{2}$ 中任意一对矢量 $\bv,\bw$，即得所需结果：

\[
\varepsilon_{ik}g^{i}_{j}v^{j}w^{k}=\varepsilon_{jm}(g^{-1})^{m}_{k}v^{j}w^{k}
\]

所以我们证明了 $SU(2)$ 的自旋-$\half$ 表示通过 $T_{\half}$ 等价于它的对偶。

由于自旋-$j$ 表示是若干份自旋-$\half$ 的对称化张量积搭建而成，我们实际上已经完成了证明 $\rho_{j}$ 也等价于其对偶所需的大部分工作。例如对 $j=1$，我们定义 $T_{1}:\mathbb{C}^{2}\otimes\mathbb{C}^{2}\to(\mathbb{C}^{2}\otimes\mathbb{C}^{2})^{*}$ 为

\[
\bigl(T_{1}(v\otimes w)\bigr)(x\otimes y)=\varepsilon_{ij}v^{i}x^{j}\,\varepsilon_{kl}w^{k}y^{l}
\]

它通过线性性延拓到 $\mathbb{C}^{2}\otimes\mathbb{C}^{2}$ 的一切元素。实际表示空间 $V_{1}$ 由 $\mathbb{C}^{2}\otimes\mathbb{C}^{2}$ 中的对称张量构成，但我们只需把 $T_{1}$ 的定义域限制为 $T_{1}:V_{1}\to V_{1}^{*}$。每个 $\varepsilon$ 在任意行列式为 1 的矩阵 $g$ 变换下不变，这一事实足以表明 $T_{1}$ 与群作用对易：

\[
T_{1}\bigl(\rho_{1}(g)\,v\otimes w\bigr)=\rho_{1}^{*}(g)T_{1}(v\otimes w)
\]

可以选取 $V_{1}$ 的一个基 $\{\bb_{1},\bb_{2},\bb_{3}\}$：

\[
\bb_{1}=\frac{\ii\,(-\be_{1}\otimes\be_{1}+\be_{2}\otimes\be_{2})}{\sqrt{2}},\quad
\bb_{2}=\frac{-(\be_{1}\otimes\be_{1}+\be_{2}\otimes\be_{2})}{\sqrt{2}},\quad
\bb_{3}=\frac{\ii\,(\be_{1}\otimes\be_{2}+\be_{2}\otimes\be_{1})}{\sqrt{2}}
\]

它按内积正交归一：

\[
\langle\bb_{i},\bb_{j}\rangle=\delta_{ij}
\]

而且若我们用 $T_{1}$ 定义一个对称双线性配对，它在另一种意义下也是"正交归一"的：

\[
\bigl(T_{1}(\bb_{i})\bigr)(\bb_{j})=\delta_{ij}
\]

所以由 $T_{1}$ 得到的矢量配对正是基 $\{\bb_{i}\}$ 下的欧几里得点积或度规。我们知道 $\rho_{1}$ 的作用是酉的，保持 $V_{1}$ 上的复内积，而 $\rho_{1}$ 还保持这个欧几里得度规：

\[
\bigl(T_{1}(\rho_{1}(g)\bb_{i})\bigr)\bigl(\rho_{1}(g)\bb_{j}\bigr)=\delta_{ij}
\]

这些基元素可以通过一个映射 $T:V_{1}\to H$ 与矩阵 $\{H_{x},H_{y},H_{z}\}$ 相联系，使得 $H_{i}=T(\bb_{i})$：

\[
T(v\otimes w)=(\sqrt{2})\,v\otimes(w^{T}\varepsilon)
\]

这里我们定义了一个分量就是列维-奇维塔符号 $\varepsilon_{ij}$ 的矩阵 $\varepsilon$，并用行矢量 $\bw^{T}$ 与 $\varepsilon$ 的矩阵乘法产生一个新行矢量。列矢量与行矢量的张量积可看成矩阵，这与"矩阵乘法作为列矢量上的线性运算"的通常观念吻合。我们可以把先前的结果

\[
\varepsilon_{ik}g^{i}_{j}=\varepsilon_{jm}(g^{-1})^{m}_{k}
\]

用矩阵乘法（而非指标记号）表达为

\[
g^{T}\varepsilon=\varepsilon g^{-1}
\]

若我们不把双覆盖同态 $f_{3}$ 描述为作用在 $\mathbb{R}^{3}$ 上，而是描述为作用在 $\{H_{x},H_{y},H_{z}\}$ 张成的矩阵空间上，则有

\[
f_{3}(g)H=gHg^{-1}
\]

而 $\rho_{1}$ 的作用是

\[
\rho_{1}(g)\,v\otimes w=gv\otimes gw
\]

我们可以用 $T$ 表明这两个表示等价：

\begin{align*}
f_{3}(g)T(v\otimes w)
&=gT(v\otimes w)g^{-1}\\
&=(\sqrt{2})\,gv\otimes(w^{T}\varepsilon)g^{-1}
=(\sqrt{2})\,gv\otimes w^{T}g^{T}\varepsilon\\
&=(\sqrt{2})\,gv\otimes(gw)^{T}\varepsilon\\
&=T(gv\otimes gw)
=T\bigl(\rho_{1}(g)\,v\otimes w\bigr)
\end{align*}

\subsection{SO(3) 与 SO(4) 的表示}

\subsubsection{SO(3) 的自旋-$j$ 表示}

对每个\emph{整数} $j$，$SU(2)$ 的玻色型自旋-$j$ 表示给出 $SO(3)$ 的一个不可约、酉的\emph{真}表示，维数为 $2j+1$。

我们在 2 对 1 函数 $f_{3}$ 的帮助下得到这个表示。对 $SO(3)$ 中任一给定旋转 $R$，都有两个互为相反数的 $SU(2)$ 元素，比如 $g$ 与 $-g$，使得 $f_{3}(g)=f_{3}(-g)=R$。然而由于 $j$ 是整数，$2j$ 为偶数，偶数个负号在 $\rho_{j}$ 的定义中相互抵消，给出

\[
\rho_{j}(-g)=(-1)^{2j}\rho_{j}(g)=\rho_{j}(g)
\]

所以选择 $g$ 还是 $-g$ 没有差别，映射 $R\mapsto\rho_{j}(g)$ 给出 $SO(3)$ 的一个维数为 $2j+1$ 的真表示。

而当 $j$ 为\emph{半奇数}时——即费米型情形——选择 $g$ 还是 $-g$ 就有差别了，因为

\[
\rho_{j}(-g)=(-1)^{2j}\rho_{j}(g)=-\rho_{j}(g)
\]

此时从 $\rho_{j}$ 只能得到 $SO(3)$ 的一个\emph{投影}的 $(2j+1)$ 维表示。

\subsubsection{SO(4) 的自旋-$(j,k)$ 表示}

$SO(4)$ 的情况非常类似。首先，对两个半整数 $j$ 与 $k$ 的某种选择，我们用 $SU(2)$ 在维数为 $2j+1$ 与 $2k+1$ 的空间 $V_{j}$、$V_{k}$ 上的自旋-$j$、自旋-$k$ 表示 $\rho_{j}$、$\rho_{k}$，在张量积 $V_{j}\otimes V_{k}$ 上定义 $SU(2)\times SU(2)$ 的一个不可约酉表示 $\rho_{j,k}$：

\[
\rho_{j,k}(g,h)\,v\otimes w=\bigl(\rho_{j}(g)v\bigr)\otimes\bigl(\rho_{k}(h)w\bigr)
\]

我们随后通过线性性把它延拓到 $V_{j}\otimes V_{k}$ 的任意元素。若 $j+k$ 为整数，该表示称为玻色型；若 $j+k$ 为半奇数，则称为费米型。

对 $SO(4)$ 中任一给定旋转 $R$，都有两对 $SU(2)$ 元素，比如 $(g,h)$ 与 $(-g,-h)$，使得 $f_{4}(g,h)=f_{4}(-g,-h)=R$。若 $j+k$ 为\emph{整数}——玻色型情形——则在这两者之间的选择无关紧要，因为

\[
\rho_{j,k}(-g,-h)=(-1)^{2(j+k)}\rho_{j,k}(g,h)=\rho_{j,k}(g,h)
\]

所以此时映射 $R\mapsto\rho_{j,k}(g,h)$ 给出 $SO(4)$ 在 $V_{j}\otimes V_{k}$ 上的一个\emph{真}表示，维数为 $(2j+1)(2k+1)$。

但当 $j+k$ 为\emph{半奇数}时——费米型情形——我们有

\[
\rho_{j,k}(-g,-h)=(-1)^{2(j+k)}\rho_{j,k}(g,h)=-\rho_{j,k}(g,h)
\]

而 $\rho_{j,k}$ 只能给出 $SO(4)$ 的一个\emph{投影}表示。

我们先前指出，从 2 对 1 同态 $f_{3}:SU(2)\to SO(3)$ 在 $\mathbb{R}^{3}$ 上得到的 $SU(2)$ 表示，等价于 $SU(2)$ 在 $V_{1}$ 上的三维自旋-1 表示。那么 2 对 1 同态 $f_{4}:SU(2)\times SU(2)\to SO(4)$ 又如何呢？结果它等价于 $SU(2)\times SU(2)$ 在 $\mathbb{C}^{2}\otimes\mathbb{C}^{2}$ 上的四维自旋-$(\half,\half)$ 表示。
